\section{Calibration of the GKT primary potential of \texorpdfstring{$x^4+y^4$}{x\^{}4+y\^{}4}}
\label{app:primary-calibration}

This appendix normalizes the source of the transport in
Section~\ref{sec:primary-frobenius}; nothing in it involves the annulus.
Appendix~\ref{subsec:reflection-fixed} shows that, for a set of labels
stable under the exchange of the coordinates, the primary parts of the GKT
potentials of the symmetric chamber indices form a single orbit under the
flows of the degree-zero field $X_{0,0}$, that this orbit contains exactly
one potential invariant under the exchange, and that these potentials define
the reflection-fixed primary potential $W^{\GKT}_{\mathrm{prim}}$ used in
Section~\ref{sec:primary-frobenius}.
Appendix~\ref{subsec:marginal-calibration} computes its restriction to the
marginal line, identifies it with the flat structure of Noumi and Yamada, and
determines the first open invariants of the reflection-fixed chamber index.
We use the notation of Section~\ref{sec:primary-frobenius}: the source-chart
coordinates $x,y$, the fields \eqref{eq:quartic-vector-field}, the primary
variables $t_{a,b}$ and the ring $A_{\mathrm{prim}}$ of
\eqref{eq:primary-base}, the good-basis cycles $\Xi_{a,b}$, the primitive-form
conditions of \cite[Def.~4.3]{grossetal2022open}, and the involution $\sigma$
and the field $X_{0,0}$ of \eqref{eq:reflection}.  The GKT notions are those
recalled in Section~\ref{subsec:gkt-recollections}.

\subsection{The reflection-fixed primary potential}
\label{subsec:reflection-fixed}

Let $I$ be a finite set of internal marking labels, each carrying a twist
$(a,b)$ with $0\leq a,b\leq2$, and let $n_{a,b}$ be the number of labels of
twist $(a,b)$.  Put
$A^{\mathrm{prim}}_I=\QQ[t_{a,b}]/(t_{a,b}^{n_{a,b}+1})$.  For a symmetric
chamber index $\nu$ bounded by $I$ \cite[Def.~3.47]{grossetal2022open}, let
$W^\nu_{\mathrm{prim}}$ be the reduction of the potential $W^{\nu,\mathrm{sym}}$
of \cite[Def.~4.12]{grossetal2022open} modulo the descendent variables
$t_{a,b,d}$, $d>0$; its coefficients lie in $A^{\mathrm{prim}}_I$.  In
\cite{grossetal2022open} \emph{symmetric} means invariance under
twist-preserving permutations of the labels
\cite[Notation~3.42]{grossetal2022open}; the exchange of the two coordinates
is the separate involution $\sigma$ of \eqref{eq:reflection}.  Let
$\mathcal P(I)$ be the set of the potentials $W^\nu_{\mathrm{prim}}$.  We
call $I$ \emph{$\sigma$-stable} if $n_{a,b}=n_{b,a}$ for all $a,b$; then
$\sigma$ acts on $A^{\mathrm{prim}}_I[x,y]$.  The $\sigma$-stable label sets
are cofinal, and $\varprojlim A^{\mathrm{prim}}_I=A_{\mathrm{prim}}$.

\begin{lemma}[Reflection-fixed representative]
\label{lem:reflection-fixed}
Let $I$ be $\sigma$-stable and put $t=t_{2,2}$.
\begin{enumerate}
\item Every element of $\mathcal P(I)$ is a polynomial of degree at most four
in $x,y$.  The automorphisms $\exp(fX_{0,0})$, with $f$ in the image of
$t^2\QQ[t^2]$ in $A^{\mathrm{prim}}_I$, act freely and transitively on
$\mathcal P(I)$.
\item The involution $\sigma$ maps $\mathcal P(I)$ to itself, and
$\sigma\exp(fX_{0,0})\sigma^{-1}=\exp(-fX_{0,0})$.
\item $\mathcal P(I)$ contains exactly one $\sigma$-fixed potential $W_I$.  If
$I\subset I'$ are $\sigma$-stable, $W_{I'}$ reduces to $W_I$.
\end{enumerate}
Consequently the $W_I$ define a polynomial
$W^{\GKT}_{\mathrm{prim}}\in A_{\mathrm{prim}}[x,y]$ of degree at most four with
$\sigma(W^{\GKT}_{\mathrm{prim}})=W^{\GKT}_{\mathrm{prim}}$, whose reduction
to every $\sigma$-stable $I$ is the primary part of a symmetric chamber index.
\end{lemma}

\begin{proof}
(1) By \cite[(0.6)]{grossetal2022open}, $W^{\nu,\mathrm{sym}}$ is homogeneous
of degree $16$ for the Euler field in which $x,y$ have weight $4$ and
$t_{a,b,0}$ has weight $16-4a-4b\geq0$.  Hence every monomial of
$W^\nu_{\mathrm{prim}}$ has degree at most four in $x,y$.

Let $\nu,\nu'$ be symmetric chamber indices bounded by $I$, and let
$W^\nu,W^{\nu'}$ be the labelled potentials of
\cite[Def.~4.6]{grossetal2022open}.  By \cite[Cor.~4.15]{grossetal2022open},
$\psi_I(W^{\nu,\mathrm{sym}})=W^\nu$, where $\psi_I$ is the injective ring
homomorphism of \cite[Lem.~4.14]{grossetal2022open}.  By
\cite[Lem.~2.5]{wuebben2026boundary}, together with
\cite[Cor.~5.14]{grossetal2022open}, the group $\widetilde G$ generated by the
wall-crossing group of \cite[Def.~4.23]{grossetal2022open} and the
bidegree-$(0,0)$ summands of \cite[Rem.~2.3]{wuebben2026boundary} acts freely
and transitively on the labelled potentials of the chamber indices bounded by
$I$.  Definition~4.23 of \cite{grossetal2022open} excludes bidegree $(0,0)$;
without these summands the group would act trivially on primary parts, as
the next paragraph shows.  Let $g\in\widetilde G$ satisfy $gW^\nu=W^{\nu'}$.
A twist-preserving permutation $\pi$ of $I$ fixes both potentials, so
$\pi(g)=g$ by freeness.

Modulo the descendent variables, a generator of $\widetilde G$ has the form
\[
 u_JX_{k_1,k_2},
 \qquad
 4k_1+4k_2=m(J,\mathbf 0)-16=\sum_{j\in J}(4a_j+4b_j-16),
\]
where $X_{k_1,k_2}$ is the field \eqref{eq:quartic-vector-field}, $u_J$ is
the product of the variables of a nonempty set $J$ of primary labels,
$k_1,k_2\geq0$, and $k_1\equiv r_J$, $k_2\equiv s_J$ modulo $4$
\cite[Def.~4.23 and Notation~2.30]{grossetal2022open}.  Here
$m(J,\mathbf d)=16+\sum_{j\in J}(4a_j+4b_j+16(d_j-1))$ is the number of
\cite[Notation~2.30]{grossetal2022open} for $r=s=4$, so these are the
conditions \eqref{eq:critical-conditions} for $\mathbf d=\mathbf 0$: the pair
$(k_1,k_2)$ is a critical exponent of $(J,\mathbf 0)$.  Since $a_j,b_j\leq2$,
the sum in the display is at most zero, with equality only if every twist in $J$ is
$(2,2)$.  Hence $(k_1,k_2)=(0,0)$, every twist in $J$ is $(2,2)$, and
$r_J=0$ forces $|J|$ to be even.  Such $J$ carry exactly one critical graph,
of critical exponent $(0,0)$ \cite[(2.31)]{grossetal2022open}, and the
corresponding generator, of bidegree $(0,0)$, is the summand $u_JX_{0,0}$
adjoined in \cite[Rem.~2.3]{wuebben2026boundary}.  These generators commute,
so modulo the descendent variables
\[
 g=\exp\Bigl(\sum_Jc_Ju_JX_{0,0}\Bigr),\qquad c_J\in\QQ.
\]
Invariance under $\pi$ gives $c_J=c_{|J|}$, and
\cite[(4.12)]{grossetal2022open} gives $\sum_{|J|=m}u_J=\psi_I(t^m/m!)$.  The
map $\psi_I$ commutes with $X_{0,0}$ and, by its monomial description
\cite[(4.12)]{grossetal2022open}, remains injective modulo the descendent
variables.  Hence
$\exp(fX_{0,0})W^\nu_{\mathrm{prim}}=W^{\nu'}_{\mathrm{prim}}$ with
$f=\sum_kc_{2k}t^{2k}/(2k)!$.

Conversely, for $f\in t^2\QQ[t^2]$ the element $\psi_I(f)X_{0,0}$ lies in the
Lie algebra of $\widetilde G$, so $\exp(\psi_I(f)X_{0,0})W^\nu=W^{\nu''}$ for a
chamber index $\nu''$ bounded by $I$.  A labelled potential determines its
chamber index, because the balanced graphs over a given marking set and
descendent vector have distinct balanced profiles.  Since
$\exp(\psi_I(f)X_{0,0})$ and $W^\nu$ are fixed by every twist-preserving
permutation, $\nu''$ is symmetric, and its primary part is
$\exp(fX_{0,0})W^\nu_{\mathrm{prim}}$.  This proves transitivity.

For freeness, let $f$ have lowest term $ct^{2k}\neq0$ in
$A^{\mathrm{prim}}_I$.  Since $W\equiv x^4+y^4$ modulo
$\mathfrak m_{\mathrm{prim}}$ and $X_{0,0}x^4=4x^4$, the coefficient of the
monomial $t^{2k}x^4$ in $\exp(fX_{0,0})W-W$ is $4c$.

(2) Write $\Gamma_{0,k_1,k_2,1,J}$, as in
\cite[Notation~2.27]{grossetal2022open}, for the smooth genus-zero graph with
$k_1$ $r$-tails, $k_2$ $s$-tails, one fully twisted boundary tail and
internal labels $J$ (Section~\ref{subsec:gkt-recollections}).  Choose a
bijection $\varsigma$ of $I$ taking each label of twist $(a,b)$ to a label of
twist $(b,a)$.  For $\nu$ symmetric define $\nu^\sigma$ by
$\nu^\sigma_{\Gamma_{0,k_1,k_2,1,J},\mathbf d}
=\nu_{\Gamma_{0,k_2,k_1,1,\varsigma(J)},\mathbf d\circ\varsigma^{-1}}$.  Since
$r=s=4$, the balancing conditions \cite[Prop.~2.31]{grossetal2022open}, the
numbers $m(J,\mathbf d)$, and the Gamma weights of
\cite[Notation~3.42]{grossetal2022open} are invariant under exchanging the two
entries of every twist and of every balanced profile.  Hence
$A(\varsigma(J),\mathbf d\circ\varsigma^{-1},\nu^\sigma)=A(J,\mathbf d,\nu)$,
the conditions of \cite[Def.~3.47]{grossetal2022open} for $\nu^\sigma$ follow
from those for $\nu$, and $\nu^\sigma$ is symmetric.  Formula
\cite[(4.11)]{grossetal2022open} gives
$W^{\nu^\sigma,\mathrm{sym}}=\sigma(W^{\nu,\mathrm{sym}})$.  The conjugation
formula follows from $\sigma(t)=t$ and $\sigma X_{0,0}\sigma^{-1}=-X_{0,0}$.

(3) Fix $W\in\mathcal P(I)$ and write $\sigma W=\exp(gX_{0,0})W$, using (1)
and (2).  For $h$ in the image of $t^2\QQ[t^2]$,
\[
 \sigma\bigl(\exp(hX_{0,0})W\bigr)
 =\exp(-hX_{0,0})\sigma W
 =\exp\bigl((g-h)X_{0,0}\bigr)W,
\]
which equals $\exp(hX_{0,0})W$ if and only if $g=2h$, by freeness.  Thus
$W_I=\exp(\tfrac12gX_{0,0})W$ is the unique $\sigma$-fixed element.  For
$I\subset I'$, restricting a chamber index to the graphs with labels in $I$
preserves the conditions of \cite[Def.~3.47]{grossetal2022open}, which for a
subset of $I$ involve only such graphs.  Restriction reduces
$W^\nu_{\mathrm{prim}}$ along $A^{\mathrm{prim}}_{I'}\to A^{\mathrm{prim}}_I$
and commutes with $\sigma$, so the image of $W_{I'}$ is a $\sigma$-fixed
element of $\mathcal P(I)$, namely $W_I$.
\end{proof}

We call $W^{\GKT}_{\mathrm{prim}}$ the \emph{reflection-fixed primary
potential}.  The variable $t_{0,0}$ has Euler weight $16$, so it occurs
only in monomials $t_{0,0}t_{2,2}^k$ of degree zero in $x,y$; write their
sum as
$c(t_{2,2})t_{0,0}$.  The $\Xi_{0,0}$ period of $W^{\GKT}_{\mathrm{prim}}$ is
$\exp(c(t_{2,2})t_{0,0}/\hbar)$ times the $\Xi_{0,0}$ period of its
restriction to $t_{0,0}=0$, and by \cite[Cor.~4.17]{grossetal2022open} the
latter is $1+O(\hbar^{-1})$.  The part of the $\hbar^{-1}$ coefficient that is
linear in $t_{0,0}$ is therefore $c(t_{2,2})t_{0,0}$, and by the same corollary
it equals $t_{0,0}$.  Hence
\begin{equation}\label{eq:constant-deformation}
 W^{\GKT}_{\mathrm{prim}}=t_{0,0}+\bigl(\text{terms free of }t_{0,0}\bigr).
\end{equation}
On the marginal line $t_{a,b}=0$, $(a,b)\neq(2,2)$, the monomials of
$W^{\GKT}_{\mathrm{prim}}$ have degree four in $x,y$, since $t_{2,2}$ has
Euler weight zero.  The congruences
$k_1\equiv r_J$, $k_2\equiv s_J$ modulo $4$ of
\cite[Prop.~2.31(a)]{grossetal2022open} allow only $t^{2k}x^4$, $t^{2k}y^4$
and $t^{2k+1}x^2y^2$, and $\sigma$-invariance equates the coefficients of
$x^4$ and $y^4$.  By \cite[Cor.~4.17]{grossetal2022open}, applied at every
finite level, $dx\wedge dy$ is a primitive form for
$W^{\GKT}_{\mathrm{prim}}$ in the sense of \cite[Def.~4.3]{grossetal2022open},
with the nine $t_{a,b}$ as flat coordinates.  By
\cite[Rem.~4.5]{grossetal2022open} these flat coordinates agree with those of
\cite[\S2.2]{heetal2022}, which come from Saito's theory of primitive forms
\cite{saito1983period}.

\begin{remark}[Reflection]
\label{rem:reflection-source}
The involution $\sigma$ fixes $W^{\GKT}_{\mathrm{prim}}$ and pulls
$dx\wedge dy$ back to $-dx\wedge dy$.  Differentiating
$\sigma(W^{\GKT}_{\mathrm{prim}})=W^{\GKT}_{\mathrm{prim}}$ gives
$\sigma(\Phi^0_{b,a})=\Phi^0_{a,b}$ for the Kodaira--Spencer classes of
Lemma~\ref{lem:source-Frobenius}.  The sign of $\sigma^*(dx\wedge dy)$ is
compensated by the exchange of $\partial_xW^{\GKT}_{\mathrm{prim}}$ and
$\partial_yW^{\GKT}_{\mathrm{prim}}$, in which the Grothendieck residue symbol
is alternating; hence the functional $\epsilon$ of
\eqref{eq:source-primary-Jacobian} satisfies
$\epsilon(\sigma h)=\sigma(\epsilon(h))$, where $\sigma$ acts on coefficients
by $t_{a,b}\mapsto t_{b,a}$.  Thus $\sigma$ is an automorphism of the source
Frobenius manifold of Lemma~\ref{lem:source-Frobenius} covering a
nontrivial involution of the base.  It is not $A_{\mathrm{prim}}$-linear; at
the Fermat point its invariants in $H_{4,4}$ form a subspace of dimension
six.
\end{remark}

\subsection{Marginal calibration}
\label{subsec:marginal-calibration}

Set $t=t_{2,2}$.  The marginal quartics are
\begin{equation}\label{eq:ABC-family}
 W_{\mathfrak a,\mathfrak b,\mathfrak a'}=\mathfrak a x^4+\mathfrak b x^2y^2+\mathfrak a'y^4,
 \qquad \Omega=dx\wedge dy.
\end{equation}
For a unit $q\in\QQ[[t]]$ the coordinate change $(x,y)\mapsto(qx,q^{-1}y)$
preserves $\Omega$ and acts by
\begin{equation}\label{eq:zero-cochain}
 (\mathfrak a,\mathfrak b,\mathfrak a')\longmapsto(q^4\mathfrak a,\mathfrak b,q^{-4}\mathfrak a').
\end{equation}
It is the monomial rescaling $z^m\mapsto q^{m_1-m_2}z^m$, a single-chart
instance of the zero-cochain action of
\cite[Prop.~5.2]{grossetal2016theta}; for $q=\exp(f)$ it is the flow
$\exp(fX_{0,0})$ of Lemma~\ref{lem:reflection-fixed}.  Its invariants are
$\mathfrak b$ and $\mathfrak a\mathfrak a'$.  An orbit with $\mathfrak a(0)=\mathfrak a'(0)=1$ has a
unique representative with $\mathfrak a=\mathfrak a'$, namely
$\mathfrak a=\mathfrak a'=\sqrt{\mathfrak a\mathfrak a'}$ with the square root equal to $1$
at $t=0$.

For the good-basis cycles $\Xi_{0,0}$ and $\Xi_{2,2}$ define the period
series
\begin{align}
 \varpi_0(\mathfrak a,\mathfrak b)
 &=\sum_{i,j,\ell\geq0}
 \frac{(-1)^{i+j}(\frac14)_{i+\ell}(\frac14)_{j+\ell}}
      {i!\,j!\,(2\ell)!}
 (\mathfrak a-1)^{i+j}\mathfrak b^{2\ell},\label{eq:F0}\\
 \varpi_2(\mathfrak a,\mathfrak b)
 &=\sum_{i,j,\ell\geq0}
 \frac{(-1)^{i+j}(\frac34)_{i+\ell}(\frac34)_{j+\ell}}
      {i!\,j!\,(2\ell+1)!}
 (\mathfrak a-1)^{i+j}\mathfrak b^{2\ell+1}.\label{eq:F2}
\end{align}
By the proof of Theorem~\ref{thm:marginal-calibration}, $\varpi_0$ and
$\varpi_2/\hbar$ are the periods of
$(W_{\mathfrak a,\mathfrak b,\mathfrak a},\Omega)$ over $\Xi_{0,0}$ and
$\Xi_{2,2}$, and the other seven periods vanish.

\begin{theorem}[Marginal period calibration]
\label{thm:marginal-calibration}
There is a unique pair $\mathfrak a(t),\mathfrak b(t)\in\QQ[[t]]$, with $\mathfrak a(0)=1$ and
$\mathfrak b(t)=t+O(t^2)$, such that
\begin{equation}\label{eq:period-inverse}
 \varpi_0\bigl(\mathfrak a(t),\mathfrak b(t)\bigr)=1,
 \qquad \varpi_2\bigl(\mathfrak a(t),\mathfrak b(t)\bigr)=t.
\end{equation}
The quartic
\[
 W_t^{\mathrm{flat}}=\mathfrak a(t)(x^4+y^4)+\mathfrak b(t)x^2y^2
\]
is the restriction of $W^{\GKT}_{\mathrm{prim}}$ to the marginal line, and
\begin{align}
 \mathfrak a(t)&=1+\frac{t^2}{16}+\frac{t^4}{768}
       +\frac{t^6}{61440}+O(t^8),\label{eq:A-series}\\
 \mathfrak b(t)&=t+\frac{t^5}{1280}+O(t^9).\label{eq:B-series}
\end{align}
For two markings of twist $(2,2)$, the open invariants of the
reflection-fixed chamber index at the balanced profiles $(4,0)$ and $(0,4)$
are
\begin{equation}\label{eq:marginal-invariants}
 \nu_{(4,0)}=\nu_{(0,4)}=-\frac18,
\end{equation}
and for three such markings its open invariant vanishes.  The first
correction visible in the Jacobian algebra of $W^{\mathrm{flat}}_t$ is the
term $t^3/32$ in the relations
\begin{equation}\label{eq:marginal-product}
 x^3=\left(-\frac t2+\frac{t^3}{32}\right)xy^2,
 \qquad
 y^3=\left(-\frac t2+\frac{t^3}{32}\right)x^2y
 \pmod{t^4}.
\end{equation}
\end{theorem}

\begin{proof}
By \cite[Lem.~4.2 and (4.3)]{grossetal2022open} the quartic good-basis
moments are
\[
 P_0(4k)=(-\hbar)^k\left(\frac14\right)_k,
 \qquad
 P_2(4k+2)=(-\hbar)^k\left(\frac34\right)_k,
\]
and vanish in the complementary congruence classes.  Expanding
\[
 \exp\!\left(
  \frac{(\mathfrak a-1)(x^4+y^4)+\mathfrak b x^2y^2}{\hbar}
 \right)
\]
against $\Xi_{0,0}$ and $\Xi_{2,2}$ gives precisely $\varpi_0$ and
$\varpi_2/\hbar$: an even number of insertions of $x^2y^2$ contributes to the
first period and an odd number to the second, and every power of $\hbar$
cancels except the displayed one.  The seven remaining periods vanish by the
phase.  Hence $dx\wedge dy$ is a primitive form for a marginal quartic with
$\mathfrak a=\mathfrak a'$, with $t$ as flat coordinate, if and only if
\eqref{eq:period-inverse} holds.

At $(\mathfrak a,\mathfrak b)=(1,0)$ the Jacobian matrix of $(\varpi_0,\varpi_2)$ with
respect to $(\mathfrak a,\mathfrak b)$ is $\operatorname{diag}(-1/2,1)$.  The formal
inverse function theorem gives existence and uniqueness; coefficient
extraction gives \eqref{eq:A-series}--\eqref{eq:B-series}.  By the
discussion after Lemma~\ref{lem:reflection-fixed}, the marginal restriction
of $W^{\GKT}_{\mathrm{prim}}$ has the form $\mathfrak a(x^4+y^4)+\mathfrak b x^2y^2$ with
$\mathfrak a(0)=1$ and $\mathfrak b=t+O(t^2)$, and it satisfies
\eqref{eq:period-inverse}; by uniqueness it is $W^{\mathrm{flat}}_t$.

At order $t^2$,
\[
 \varpi_0=1-\frac12(\mathfrak a-1)+\frac{\mathfrak b^2}{32}+O(t^3),
\]
so $\mathfrak a-1=t^2/16+O(t^3)$.  By the sign and the order-two automorphism
factor in \cite[(4.11)]{grossetal2022open}, the coefficient of $t^2x^4$,
respectively $t^2y^4$, is $-1/2$ times the corresponding invariant for two
markings.  Hence \eqref{eq:marginal-invariants}.  The coefficient of $t^3$
in the second period equation gives the vanishing of the invariant for three
markings.  Finally the Jacobian relation is $x^3=-\mathfrak b(2\mathfrak a)^{-1}xy^2$, and expansion
of $-\mathfrak b/(2\mathfrak a)$ gives \eqref{eq:marginal-product}.
\end{proof}

In the coordinates $\tilde x=\mathfrak a(t)^{1/4}x$, $\tilde y=\mathfrak a(t)^{1/4}y$
the pair $(W^{\mathrm{flat}}_t,dx\wedge dy)$ becomes
$(\tilde x^4+\tilde y^4+s\,\tilde x^2\tilde y^2,\;
\mathfrak a(t)^{-1/2}d\tilde x\wedge d\tilde y)$ with $s=\mathfrak b(t)/\mathfrak a(t)$.  The
same substitution in the moments gives
$\varpi_0(\mathfrak a,\mathfrak b)=\mathfrak a^{-1/2}\varpi_0(1,s)$ and
$\varpi_2(\mathfrak a,\mathfrak b)=\mathfrak a^{-1/2}\varpi_2(1,s)$, while
\[
 \varpi_0(1,s)={}_2F_1\!\left(\tfrac14,\tfrac14;\tfrac12;\tfrac{s^2}4\right),
 \qquad
 \varpi_2(1,s)=s\,{}_2F_1\!\left(\tfrac34,\tfrac34;\tfrac32;\tfrac{s^2}4\right).
\]
Hence \eqref{eq:period-inverse} is equivalent to
\begin{equation}\label{eq:NY-form}
 \mathfrak a(t)^{1/2}={}_2F_1\!\left(\tfrac14,\tfrac14;\tfrac12;\tfrac{s^2}4\right),
 \qquad
 t=\frac{s\,{}_2F_1\bigl(\frac34,\frac34;\frac32;\frac{s^2}4\bigr)}
        {{}_2F_1\bigl(\frac14,\frac14;\frac12;\frac{s^2}4\bigr)} .
\end{equation}
The second relation is the flat coordinate of the simple elliptic
singularity $x^4+y^4$ along its marginal deformation computed by Noumi and
Yamada \cite{noumiyamada1998flat}, as quoted in \cite[(7)]{strachan2010}.
The statements of \cite[Props.~4.2.6--4.2.8]{maher2024predictions} for
elliptic singularities of Fermat type (the primitive form as the reciprocal
of a Gauss hypergeometric function times the volume form after Morse
stabilization, the flat coordinate, and $s(t)$ as a Schwarz triangle
function) specialize here to ${}_2F_1(\frac14,\frac14;\frac12;z)$,
$z=s^2/4$.  By the first relation in \eqref{eq:NY-form}, the resulting
primitive form
$d\tilde x\wedge d\tilde y/{}_2F_1(\frac14,\frac14;\frac12;\frac{s^2}4)$ is
$dx\wedge dy$, the normalization of \cite{grossetal2022open} used in
Theorem~\ref{thm:marginal-calibration}.  In this normalization the
coefficients of the potential are determined by the open invariants of the
reflection-fixed chamber index through \cite[(4.11)]{grossetal2022open}, and
the series \eqref{eq:A-series}--\eqref{eq:B-series} and the invariants
\eqref{eq:marginal-invariants} refer to it.

The residue pairing detects this normalization at quadratic order.  For
$W_c=(1+ct^2)(x^4+y^4)+tx^2y^2$, the normalized residue of the socle
monomial is
\begin{equation}\label{eq:residue-gauge}
 16\operatorname{Res}_0
 \frac{x^2y^2\,dx\wedge dy}{\partial_xW_c\,\partial_yW_c}
 =\frac1{(1+ct^2)^2-t^2/4}.
\end{equation}
Since $x^4\equiv-\frac{t}{2(1+ct^2)}x^2y^2$ in the Jacobian algebra, the
residue pairing $\eta$ of $(W_c,dx\wedge dy)$ satisfies
\begin{equation}\label{eq:KS-gauge}
 \eta(\partial_tW_c,1)
 =\frac{1-2ct^2/(1+ct^2)}{(1+ct^2)^2-t^2/4}
 =1+\left(\frac14-4c\right)t^2+O(t^4).
\end{equation}
When $dx\wedge dy$ is a primitive form and $t$ a flat coordinate, this entry of
the flat metric is constant, equal to its value $1$ at $t=0$.  Its constancy
forces $c=1/16$ at order $t^2$, whereas the product first sees $c$ at order
$t^3$: in \eqref{eq:marginal-product} the relevant coefficient is
$-t/(2(1+ct^2))=-t/2+(c/2)t^3+O(t^5)$.

Off the marginal line the flat coordinates are simultaneous, and on the first
two-parameter slice this already fixes a term that the critical locus does
not see.  Set $u=t_{1,1}$ and $t=t_{2,2}$.  Through total degree three, the
restriction of $W^{\GKT}_{\mathrm{prim}}$ to the slice on which the other
seven $t_{a,b}$ vanish is
\begin{equation}\label{eq:mixed-primary-potential}
 W_{u,t}
 =\left(1+\frac{t^2}{16}\right)(x^4+y^4)
  +tx^2y^2
  +\left(u-\frac{ut^2}{16}\right)xy
  -\frac{u^2t}{32}+O((u,t)^4).
\end{equation}
The two new open invariants of the reflection-fixed chamber index are
\[
 \nu_{ut^2;(1,1)}=-\frac18,
 \qquad
 \nu_{u^2t;(0,0)}=-\frac1{16}.
\]
They are the solutions of the chamber conditions
\cite[Def.~3.47(3)]{grossetal2022open} for the supports $J$ consisting of
one marking of twist $(1,1)$ and two of twist $(2,2)$, respectively two of
twist $(1,1)$ and one of twist $(2,2)$; both supports satisfy the
inequality in that condition.  These supports have the single balanced
profiles $(1,1)$ and $(0,0)$, and no two-element support containing a
marking of twist $(1,1)$ has a balanced profile.  In the expression
$A(J,\mathbf 0,\nu)$ of Section~\ref{subsec:gkt-recollections}, with the
weights \eqref{eq:quartic-Gamma-weight}, the two conditions therefore read
\[
 \nu_{ut^2;(1,1)}+\tfrac12\bigl(\nu_{(4,0)}+\nu_{(0,4)}\bigr)+\tfrac14=0,
 \qquad
 \nu_{u^2t;(0,0)}+\tfrac1{16}=0,
\]
where the middle term comes from the partition that separates the marking
of twist $(1,1)$, the terms $\frac14$ and $\frac1{16}$ come from the
partitions into singletons, and $\nu_{(4,0)}=\nu_{(0,4)}=-\frac18$ by
\eqref{eq:marginal-invariants}.  The repeated-marking automorphisms in
\cite[(4.11)]{grossetal2022open} divide the two invariants by two, giving
the last two corrections in \eqref{eq:mixed-primary-potential}.  The
periods confirm these values.  For the $\Xi_{1,1}$ period, the direct
correction $-\frac{ut^2}{16}xy$, the cross term of $uxy$ with
$\frac{t^2}{16}(x^4+y^4)$, and the cubic term $uxy(tx^2y^2)^2$ contribute
\[
 -\frac{ut^2}{16\hbar},\qquad
 -\frac{ut^2}{16\hbar},\qquad
 \frac{ut^2}{8\hbar};
\]
for $\Xi_{0,0}$ the two contributions are
$-u^2t/(32\hbar)$ and $u^2t/(32\hbar)$.  Both sums vanish, as the flatness
of $u$ and $t$ requires.  Thus the constant term $-u^2t/32$ is fixed by the
oscillatory integral although it is invisible to the critical locus.
